\section{Introduction}\label{sec:introduction}

The two-dimensional nonlinear sigma model asks how a field constrained to
a curved compact target can produce a relativistic quantum field theory
at distances much larger than a lattice spacing. For the sphere $S^2$, the
basic lattice model has a particularly simple definition. On a finite
periodic square lattice $\Lambda$, let $q_x\in S^2\subset\R^3$ and let
$\sigma$ be normalized area measure on $S^2$. Its Gibbs probability is
\begin{equation}\label{eq:lattice-model}
 \dd\mu_{\beta,\Lambda}(q)
 =\frac{1}{Z_{\beta,\Lambda}}
   \exp\!\left(\beta\sum_{\{x,y\}\in E(\Lambda)}q_x\cdot q_y\right)
   \prod_{x\in\Lambda}\dd\sigma(q_x),
 \qquad \beta>0,
\end{equation}
where each nearest-neighbor bond occurs once. We consider no topological
term and no external field. The continuum problem is to let the coupling
$\beta$ diverge and the lattice spacing vanish while preserving a
nontrivial physical scale. A mass estimate at fixed lattice spacing does
not settle this problem: one must construct the same limiting fields on
which the limiting gap acts.

\subsection{Historical context}
\label{subsec:history}

The nonordering theorem of Mermin and Wagner, and Mermin's classical
version, exclude spontaneous magnetization in the corresponding
finite-range two-dimensional systems~\cite{MW,Mermin1967}.
McBryan and Spencer obtained algebraic upper bounds on spin correlations
by complex rotations~\cite{McBryanSpencer}.  Such an upper bound does not
distinguish algebraic decay from exponential decay.  Local Ward
identities gave further correlation inequalities and high-temperature
mass-gap criteria in the work of Aizenman and
Simon~\cite{AizenmanSimonWard}.  A continuum mass theorem additionally
requires control as the bare inverse coupling diverges and the lattice
spacing vanishes.

For more than two spin components, perturbative renormalization
predicts that weak microscopic coupling produces an exponentially
large correlation length.  Polyakov identified the renormalization-group
mechanism for two-dimensional Goldstone fields~\cite{Polyakov}.
Br\'ezin and Zinn-Justin developed the expansion near two dimensions,
and Br\'ezin, Zinn-Justin and Le Guillou analyzed the renormalization of
the model and its composite operators~\cite{BrezinZinnJustin,BrezinZinnJustinLeGuillou}.
These works explain why weak microscopic coupling can coexist with a
finite long-distance mass scale.  They also identify the perturbative
coefficients against which a constructive renormalization should be
checked.  Establishing those coefficients is logically different from
bounding the nonperturbative remainder throughout a trajectory whose
length diverges with the cutoff.

Integrability gives a complementary and much more detailed picture of
the expected continuum theory.  Zamolodchikov and Zamolodchikov solved
the factorization, crossing and unitarity equations for
$O(N)$-invariant scattering and identified the sigma-model
solution~\cite{ZamolodchikovZamolodchikov}.  Hasenfratz, Maggiore and
Niedermayer obtained the exact mass-to-renormalization-scale relation
for $O(3)$ and $O(4)$ by matching the Bethe ansatz with perturbation
theory~\cite{HasenfratzMaggioreNiedermayer}; Hasenfratz and Niedermayer
extended the calculation to arbitrary $N\ge3$~\cite{HasenfratzNiedermayer}.
Those calculations determine a mass ratio within the integrable
continuum description.  The additional problem addressed here is to
construct the field distributions from the specified nearest-neighbor
measure and to establish their spectral properties after removal of
the cutoff.  No factorized scattering matrix is used as a premise of
the construction.

Susceptibility and second-moment correlation length provide intrinsic
field and length units.  This is the familiar zero-momentum normalization,
used explicitly by Campostrini, Pelissetto, Rossi and
Vicari~\cite{CampostriniPelissettoRossiVicari1997}.
Their strong-coupling calculations and the lattice--bootstrap comparisons
of Balog and collaborators~\cite{BalogEtAl1999} studied normalized
amplitudes and supplied evidence for continuum universality.
Here the question is convergence of the entire canonically normalized
Schwinger hierarchy for the specified nearest-neighbor model and periodic
thermodynamic state, as the bare coupling tends to infinity without
restriction to a selected sequence.  Comparing this limit with a
terminal-coupling construction requires control of both the field
distributions and the normalizing correlation moments.

Rigorous work has supplied several distinct parts of this program.
Kupiainen proved asymptotic validity of the $1/N$ expansion for lattice
models above the spherical-model critical temperature and obtained a
mass gap for sufficiently large $N$ at each such
temperature~\cite{Kupiainen}.
Kopper later proved mass generation at sufficiently large finite $N$
with a suitable ultraviolet cutoff~\cite{Kopper}.  These large-$N$
results do not specialize to a cutoff-removed construction at $N=3$.
Gaw\k{e}dzki and Kupiainen constructed a continuum limit with asymptotic
freedom for a sigma model with a \emph{hierarchical} kinetic
term~\cite{GawedzkiKupiainen}.  In that setting the hierarchical covariance reduces the flow to
a single-spin recursion.  For the nearest-neighbor model, integration
generates interactions between blocks whose spatial decay must be
controlled.
Mitter and Ramadas constructed the Wilson renormalized trajectory in
perturbation theory in the effective charge and developed the
corresponding effective-action analysis~\cite{MitterRamadas}.

The organization by successive block integrations belongs to the
renormalization-group tradition of Kadanoff and
Wilson~\cite{Kadanoff,Wilson}.  Ba\l aban developed localized effective actions and constrained
variational methods in lattice gauge theory~\cite{BalabanGaugeFlow,BalabanGaugeVariational}
and in the low-temperature analysis of classical $N$-vector
models~\cite{BalabanFlow,BalabanVariational}.
Goswami studied the constrained background-field problem for
two-dimensional $SU(n)$ chiral models, including the $O(4)$
case~\cite{Goswami2024}.
For the two-dimensional $O(4)$ model, Dybalski, Stottmeister and
Tanimoto proved existence and uniqueness for the constrained
small-field variational problem associated with one block-averaging
operation~\cite{DybalskiStottmeisterTanimotoVariational}.
Aru, Garban and Sep\'ulveda proved that, after fixing one spin
direction, rescaled transverse fluctuations converge locally as the
inverse coupling tends to infinity to a rooted vector-valued lattice
Gaussian free field~\cite[Theorem~1.3]{AGS}.  For $O(3)$ this field has
two components.  Their result identifies the local Gaussian
approximation; the construction below also requires control at
spatial scales that grow with the inverse coupling.


\subsection{Canonical normalization and the main result}

For each $\beta$ used below, write $\mu_\beta$ for the local limit of
periodic square tori as their side lengths tend to infinity. Existence and
uniqueness of this periodic limit are part of the preliminary estimates.
Given a lattice spacing $a>0$, a positive field normalization $B$, and
$f\in C_c^\infty(\R^2)$, define the three-component random variable
\[
 \phi_{a,B}(f)=Ba^2\sum_{x\in\Z^2}f(ax)q_x.
\]
For component indices $i_1,\ldots,i_n\in\{1,2,3\}$, the corresponding
Schwinger distribution is the expectation of the product of these
smeared fields. Its extension to general tests on $(\R^2)^n$ is equivalently
the discrete moment measure with weight $(Ba^2)^n$.

The lattice itself supplies two normalization constants.  Put
\begin{equation}\label{eq:canonical-lattice-scales}
 C_\beta(x)=\E_{\mu_\beta}[q_0^1q_x^1],\qquad
 \chi_\beta=\sum_{x\in\Z^2}C_\beta(x),\qquad
 \xi_\beta^2=\frac{1}{4\chi_\beta}
                 \sum_{x\in\Z^2}|x|^2 C_\beta(x).
\end{equation}
Here $\chi_\beta$ is the susceptibility of one spin component, and
$\xi_\beta$ is the positive second-moment correlation length.  For a
vector test $f\in C_c^\infty(\R^2;\R^3)$ define
\begin{equation}\label{eq:canonical-field}
 \Psi_\beta(f)=\frac{1}{\xi_\beta\sqrt{\chi_\beta}}
           \sum_{x\in\Z^2} f(x/\xi_\beta)\cdot q_x.
\end{equation}
The factor $4$ in \eqref{eq:canonical-lattice-scales} is twice the Euclidean
dimension.  For a translation-invariant continuum two-point measure
written in relative coordinates as $\dd y\,\nu(\dd z)$, with
$z=y_2-y_1$, its susceptibility and squared second-moment length are
$\nu(\R^2)$ and $\int|z|^2\nu(\dd z)/(4\nu(\R^2))$, respectively,
whenever these quantities are finite and positive.

\begin{theorem}[Canonical continuum limit]\label{thm:canonical}
For the nearest-neighbor probability \eqref{eq:lattice-model}, the
periodic thermodynamic limit $\mu_\beta$ exists for all sufficiently
large $\beta$.  Both sums in \eqref{eq:canonical-lattice-scales} are
finite, $\chi_\beta,\xi_\beta>0$, and $\xi_\beta\to\infty$ as
$\beta\to\infty$.  The following limits hold through all real values
of $\beta$ tending to infinity.
\begin{enumerate}[label=\textup{(\roman*)}]
\item For every finite list of vector tests in
$C_c^\infty(\R^2;\R^3)$, the fields \eqref{eq:canonical-field} converge
jointly in law and in all joint moments to one common limiting hierarchy.
\item For every order and every choice of spin components, the associated
Schwinger distributions converge in the strong topology of
$\mathcal S'((\R^2)^n)$, that is, uniformly on bounded sets of Schwartz
tests.
\item The limiting hierarchy has susceptibility and second-moment length
equal to one.  It is the hierarchy constructed in Theorem~\ref{thm:main},
after the unique positive constant length and field rescalings imposing
these two normalizations.
\end{enumerate}
\end{theorem}

The assertion concerns the one lattice action and periodic thermodynamic
state specified above.  The normalization removes the freedom to choose
constant physical length and field units.  It does not introduce a
comparison with other lattice actions, external fields, or topological
terms.

\subsection{Construction and physical properties}

The proof first constructs the fields at a prescribed physical block
scale.  This gives the following more detailed existence statement,
including the physical properties that will pass to the canonical limit.

\begin{theorem}\label{thm:main}
There are a dyadic integer $L>1$, a constant $H>0$, bare couplings
$\beta_N\to\infty$, and field normalizations $B_N>0$, such that
\begin{equation}\label{eq:main-prescription}
 a_N=L^{-N},\qquad
 \beta_N=H+\frac{\log L}{2\pi}N+O(\log(N+1)),
\end{equation}
and the following assertions hold for the model
\eqref{eq:lattice-model}.
\begin{enumerate}[label=\textup{(\roman*)}]
\item Every finite list of the variables
$\phi_N(f)=\phi_{a_N,B_N}(f)$, under $\mu_{\beta_N}$, converges jointly
in law and in all moments. The component Schwinger distributions converge
in $\mathcal S'((\R^2)^n)$ at every order. The same limits are obtained
with suitable periodic physical tori tending to the plane as $N\to\infty$.
\item The limiting Schwinger distributions are Euclidean invariant,
internally $O(3)$ invariant, reflection positive, symmetric and clustering.
For every component distribution and $F\in\mathcal S((\R^2)^n)$,
\[
 |S_n(F)|\le C^n n!\sup_{x\in(\R^2)^n}
       (1+|x|^2)^{2n}|F(x)|,
\]
with a constant $C$ independent of $n$. They reconstruct a local, unitary,
relativistic theory in $1+1$ dimensions.
\item The reconstructed theory has a unique vacuum $\Omega$, a nonzero
vacuum complement, and a Hamiltonian $H_{\mathrm{phys}}$ satisfying
\[
 H_{\mathrm{phys}}\big|_{\Omega^\perp}\ge m\Id
 \quad\text{for some }m>0.
\]
\item A connected four-point distribution is nonzero on smooth tests
with strictly separated time supports. In particular, the limiting
field correlations do not satisfy Wick factorization.
\end{enumerate}
The trajectory is selected by fixing its terminal kinetic coupling to
$H$ at a fixed physical block length.
\end{theorem}

Theorem~\ref{thm:main} constructs a continuum limit along a prescribed
trajectory.  Theorem~\ref{thm:canonical} identifies its canonical
normalization with the limit along every diverging bare coupling.
Interaction is expressed by the fourth assertion of
Theorem~\ref{thm:main}: a connected four-point correlation of the
constructed field is nonzero.  The mass-gap estimate and this criterion
use the same fields and the same choice of physical units.

\begin{corollary}\label{cor:canonical-physics}
The canonical limiting hierarchy is Euclidean invariant, internally
$O(3)$ invariant, reflection positive, symmetric, and clustering.  It
satisfies factorial tempered-distribution bounds and reconstructs a local,
unitary relativistic theory in $1+1$ dimensions.  The reconstructed
Hamiltonian has a unique vacuum, a nonzero vacuum complement, and a
strictly positive lower bound on that entire complement.  A connected
four-point distribution is nonzero on smooth tests with strictly separated
time supports.
\end{corollary}

\subsection{Proof strategy and new estimates}

The starting analytic tools are the local angular-integration and exact
blocking constructions in \cite{OpenAI-O4}, denoted by R when discussing
their adaptation. We use their stated intermediate estimates with the
hypotheses specified below. We do not transfer the $O(4)$ mass theorem to
$O(3)$. The difference in target dimension changes the current identities,
the Gaussian normalization and the coefficient of the coupling drift.
Sections~\ref{sec:preliminary} and~\ref{sec:rg} establish these changes
directly. Section~\ref{sec:setup} fixes the imported definitions and
identifies the dimension-independent inputs.

An exact blocking partitions the lattice into squares, introduces one
retained spin in $S^2$ per square through a normalized averaging kernel,
and integrates the fine spins. This preserves the partition function and,
when sources are retained, the generating function of the original spins.
One carries those sources and all the corrections they generate through
successive integrations. At a step with side factor $l$, rescale the
source at each fine site in a block to $l^{-2}$ times its retained source
$z$. On a constant retained-spin configuration $V$, the resulting
linear term is $c_{j,N}z\cdot V$. Normalize its coefficient to one.
Section~\ref{sec:sources} constructs these positive coefficients and
proves that the accumulated field normalization is
$B_N=\prod_{j=1}^N c_{j,N}^{-1}$.

The proof must connect a rough large-distance estimate with a precise
comparison at a fixed physical scale. Four additional estimates make that
connection possible.

First, write $X(\beta)=C\exp((4\pi-\eta)\beta)$ for a preliminary
upper bound on the mixing length, where $\eta>0$ is fixed.
The strict saving below $4\pi$ comes from the transverse fluctuations in
the first coarse tile. The pinned-cell inputs are uniform over boundary
pins. The resulting mixing estimate is uniform under cuts and weaker bonds.

Second, spatial inversion and tangent reflection give contraction of
the terms remaining after extraction of the kinetic coupling and its
specified low-order corrections. The contraction can be made arbitrarily
close to $L^{-2}$ per step. The
affine quadratic contribution has already been removed by the kinetic
normalization; inversion cancels the next spatial cross term. This
additional cancellation allows the coarse comparison box to grow while
the total error in its density still tends to zero.

Third, an integrated comparison controls changes in the gas of
large-gradient regions. We fill selected regions with small-gradient
spins, retain the complete positive density at the filled configuration,
and pay the filling entropy with the original gradient energy. The
comparison loses only a volume factor. The strict saving in the
preliminary length permits a reference box larger than that mixing length,
while the nearly $L^{-2}$ contraction keeps the error over its coarse
volume small. Together these estimates transfer a finite-box criterion
for exponential decay from one reference cutoff to every finer cutoff
at one fixed physical box. The transfer-operator criterion then gives a
common physical decay rate. More concretely, at a reference depth $K$,
the preliminary mixing length in physical units is
\[
 a_KX(\beta_K)=L^{(1-\eta/(2\pi))K+o(K)}.
\]
Matching the effective densities to finer cutoffs is controlled by the
error parameter $C_H M^2L^{-(2-\upsilon)K}$ in a physical square of side
$M$, where $\upsilon>0$ can be chosen arbitrarily small. Thus $M$ can be
much larger than the mixing length while this error tends to zero.
Fixing one sufficiently large $K$ fixes a successful physical square
once and for all; its decay criterion passes to every finer cutoff.

Fourth, independent local sources are carried through the exact
transformation. Their linear response defines $B_N$; the remaining source
terms contract. This gives convergence of the actual smeared lattice
spins, local exponential moments, and a fourth cumulant that survives the
limit. Two reflected inclined first blockings compare regulator
orientations through a common coarse frame. The resulting rotation has
angle equal to an irrational multiple of $\pi$. Together with continuity,
it gives full Euclidean rotation invariance.

The resulting architecture separates the long-distance estimate from
the source construction. Sections~\ref{sec:free}--\ref{sec:shooting}
construct and match the exact trajectories.
Sections~\ref{sec:comparison}--\ref{sec:trace} obtain the estimate uniform
in the cutoff. Sections~\ref{sec:sources}--\ref{sec:continuum} construct
the Schwinger distributions and prove their symmetries.
Section~\ref{sec:os} reconstructs the quantum theory and proves both the
full spectral gap and non-Wick factorization.

The final step is to make the physical normalization intrinsic to the
lattice.  An arbitrary large bare coupling can be written as a reference
coupling at an integer depth plus a bounded real offset.  The effective
hierarchy initially may depend on that offset.  A summable sensitivity
estimate for the noncanonical kinetic increment, proved in
Lemma~\ref{lem:sensitivity}, shows that trajectories with the same
corrected offset agree at each fixed coarse scale.  This is stronger
than a bound on the size of the increment: it controls its change under
perturbations of the trajectory.  The resulting argument in
Section~\ref{sec:trajectories} is useful whenever a marginal coupling
must be matched over arbitrarily many renormalization steps.

Section~\ref{sec:volume} extends the fixed-scale volume comparison to
these bounded offsets.  Section~\ref{sec:fields} then constructs their
continuum hierarchies and proves uniform two-point tails.  The tails
permit passage of the full susceptibility and second moment to the
limit, so the normalization in \eqref{eq:canonical-field} can be applied
after taking the continuum limit.  Finally,
Section~\ref{sec:uniqueness} compares two descriptions of the same
microscopic field: adjacent ordinary depths, and an ordinary versus an
inclined first blocking.  Besides rotating the lattice, the inclined
step contributes a length factor $5$, multiplicatively independent of
the dyadic factor $L$.  The two reflected inclinations remove the
rotation; continuity and the two resulting incommensurable offset
periods remove the remaining offset dependence.  Thus the same exact
blocking geometry serves both Euclidean symmetry and canonical
uniqueness.

The conditional comparisons in the preliminary mixing argument use
Ginibre's formulation of ferromagnetic correlation inequalities and
the positive-association theorem of Fortuin, Kasteleyn and Ginibre~\cite{Ginibre1970,FortuinKasteleynGinibre1971}.
The passage from local disagreement estimates to loss of boundary
influence is related to disagreement percolation~\cite{VanDenBergMaes1994}.
In exact blocking, the connected polymer sums use the classical
tree-graph approach to cluster expansions; a useful formulation of
the Penrose identity is given by Fern\'andez and
Procacci~\cite{FernandezProcacci}.

The finite-volume comparison uses reflection positivity and the
chessboard method systematized by Fr\"ohlich, Israel, Lieb and
Simon~\cite{FILS78,FILS80}.  For the final passage from Euclidean
distributions to a local relativistic field, we use the reconstruction
framework of Osterwalder and Schrader~\cite{OS73} in the corrected
form with the linear growth condition of \cite[Section~IV.1]{OS75}.
Reflection positivity, Euclidean covariance, clustering and the
hierarchy of distributional bounds are consequently separate parts
of the continuum argument.

